% Optimizer reference and synthetic-benchmark documentation for the
% geko_bayopt project.
%
% Written as a standalone, appendix-style document (top-level headings are
% \section, not \chapter) so it can be compiled on its own (e.g. via Overleaf)
% or have its body \input{} into an existing report.  If embedding, demote
% every \section -> \subsection as needed and drop the preamble down to
% \begin{document}.
%
% Figures are read from optimizer_visualization/plots/individual/ --
% compile from within optimizer_visualization/, or adjust \graphicspath.
%
% Companion document: optimizer_comparison_report.tex (real-CFD comparison).
% This document covers (1) the configuration reference for every optimizer in
% src/geko_bayesopt/optimizer.py and (2) the synthetic benchmark implemented
% in _benchmark_core.py / benchmark_individual.py.

\documentclass[11pt]{article}

\usepackage[margin=1in]{geometry}
\usepackage{graphicx}
\usepackage{booktabs}
\usepackage{amsmath}
\usepackage{listings}
\usepackage{xcolor}
\usepackage{hyperref}
\hypersetup{colorlinks=true, linkcolor=blue, citecolor=blue, urlcolor=blue}

\graphicspath{{plots/individual/}}

\lstset{
  basicstyle=\ttfamily\small,
  columns=fullflexible,
  breaklines=true,
  frame=single,
  framerule=0.3pt,
  rulecolor=\color{black!30},
  xleftmargin=1em,
  aboveskip=6pt,
  belowskip=6pt,
}

\newcommand{\csep}{C_{sep}}
\newcommand{\cnw}{C_{nw}}
\newcommand{\cmix}{C_{mix}}
\newcommand{\cjet}{C_{jet}}
\newcommand{\ccorner}{C_{corner}}
\newcommand{\cturb}{C_{turb}}

\title{Optimizers for GEKO Coefficient Calibration\\[2pt]
\large Configuration Reference and Synthetic Benchmark}
\author{}
\date{}

\begin{document}
\maketitle
\tableofcontents

% ======================================================================
\section{Overview and Shared Architecture}
\label{sec:architecture}

The package \texttt{geko\_bayesopt} provides eight optimizer kinds, all
implemented in \path{src/geko_bayesopt/optimizer.py}: four standalone
strategies (Bayesian Optimization, Nelder--Mead, finite-difference gradient
descent, Particle Swarm Optimization) and four two-phase hybrids that combine
Bayesian Optimization with one of the heuristic methods in either order.
This section documents what all of them share; Section~\ref{sec:optimizers}
documents each kind individually, and
Sections~\ref{sec:benchmark}--\ref{sec:benchmark-results} document the
synthetic benchmark used to validate them.

\subsection{The ask/tell interface}

Every optimizer implements the same four-method protocol:

\begin{lstlisting}
x = opt.ask()          # suggest the next parameter vector
opt.tell(x, y)         # report the objective value y observed at x
opt.should_stop()      # True when the epsilon convergence check fires
opt.test_config()      # raises ValueError on an invalid configuration
\end{lstlisting}

The experiment loop strictly alternates
$\texttt{ask()} \rightarrow \text{evaluate} \rightarrow \texttt{tell()}$;
none of the optimizers support batched or parallel asks.
\texttt{build\_optimizer(section, parameters)} constructs the requested kind
from the JSON configuration and calls \texttt{test\_config()} before
returning, so invalid configurations fail at build time rather than
mid-run.

A further design invariant worth knowing about (it explains several
implementation choices below): interrupted experiments are resumed by
\emph{replaying} every completed trial through \texttt{tell()} on a fresh
optimizer object; \texttt{ask()} is never called during a replay.  Each
optimizer therefore advances its internal state machine entirely inside
\texttt{tell()}, and all eight kinds resume bit-for-bit identically to an
uninterrupted run.

\subsection{Configuration shape}

All optimizers are configured inside the \texttt{"optimizer"} block of an
experiment JSON, which always has the same top-level shape:

\begin{lstlisting}
"optimizer": {
    "kind": "<optimizer_kind>",
    "stopping_criteria": {
        "n_calls": 32,
        "epsilon": 1e-4,
        "window": 3
    },
    "kind_specific_options": { ... }
}
\end{lstlisting}

\texttt{stopping\_criteria} and \texttt{kind\_specific\_options} may be
omitted entirely to use the defaults.  The valid \texttt{kind} values are
\texttt{skopt\_gp}, \texttt{nelder\_mead}, \texttt{finite\_differences},
\texttt{pso}, \texttt{hybrid\_nm\_bayes}, \texttt{hybrid\_fd\_bayes},
\texttt{hybrid\_bayes\_nm}, and \texttt{hybrid\_bayes\_fd}.

\subsection{Stopping criteria}
\label{sec:stopping}

\begin{center}
\begin{tabular}{@{}llp{9cm}@{}}
\toprule
Field & Default & Description \\
\midrule
\texttt{n\_calls} & 32 & Hard cap on total objective (CFD) evaluations. The
experiment always stops here at the latest. \\
\texttt{epsilon} & $10^{-4}$ & Relative-improvement threshold for early
stopping. Must be $>0$; set to \texttt{null} to disable. \\
\texttt{window} & 3 & Number of \emph{meaningful steps} compared by the
epsilon check. Must be $\geq 1$. \\
\bottomrule
\end{tabular}
\end{center}

After each trial the optimizer checks whether improvement has stalled: the
best value within the last \texttt{window} meaningful steps is compared
against the best value before those steps, and if the relative improvement
is below \texttt{epsilon} the experiment stops early.  At least
$2\times\texttt{window}$ meaningful steps must exist before the check can
fire, which prevents spurious stops at the very beginning.  Because the
comparison uses the absolute difference of the two best values, a
\emph{worsening} recent window (possible for finite differences, whose step
series is not monotonic) does not trigger a stop; only stagnation does.

What counts as a meaningful step differs by optimizer, because raw function
evaluations are not always decisions:

\begin{center}
\begin{tabular}{@{}lp{10.5cm}@{}}
\toprule
Optimizer & One meaningful step $=$ \\
\midrule
\texttt{skopt\_gp} & every evaluation \\
\texttt{nelder\_mead} & one simplex iteration (startup evaluations and
individual shrink evaluations do not count) \\
\texttt{finite\_differences} & one complete gradient cycle ($D$ probes $+$ 1
step); probe evaluations are invisible to the check \\
\texttt{pso} & one completed swarm iteration (the global best after all
\texttt{n\_particles} results are in) \\
hybrids & the active phase's rule above \\
\bottomrule
\end{tabular}
\end{center}

\subsection{Boundary policy}
\label{sec:bounds}

The parameter bounds declared in the \texttt{parameters} block are handled
differently per optimizer, deliberately:

\begin{itemize}
  \item \textbf{Bayesian Optimization} only ever proposes points inside the
        bounds, by construction of the search space.
  \item \textbf{Nelder--Mead and finite differences do \emph{not} clip}:
        a reflection, expansion, probe, or gradient step is evaluated exactly
        where proposed, even outside the nominal range.  A poor score there
        steers the search back naturally.  Only the \emph{startup} points
        (initial simplex, initial FD base) are clipped into the bounds so the
        search begins in the valid region.
  \item \textbf{PSO absorbs}: a particle that overshoots a bound is placed
        exactly on it and its velocity component in that dimension is zeroed.
        This is part of the PSO algorithm itself, not a harness workaround.
\end{itemize}

The practical consequence: with NM or FD (standalone or inside a hybrid),
the search can in principle wander far outside the declared bounds.  The
only hard limit is the point at which Fluent no longer accepts the proposed
coefficient values; this has not happened in practice yet.  The best-found
point may accordingly also lie outside the declared bounds, if that is
where the best score occurred.

\subsection{Determinism}

All optimizers are bitwise reproducible given the same configuration: every
stochastic component (the Sobol' initial samples, the GP
hyperparameter-fit restarts, PSO's velocity draws) is seeded from the
configuration's \texttt{random\_state}.  In addition, all five GP
construction sites (standalone \texttt{skopt\_gp} and the four hybrids)
build the identical estimator (Mat\'ern $\nu=2.5$ kernel, 10 restarts of
the kernel hyperparameter fit), so every optimizer's BO component behaves
the same and identical configurations produce identical trajectories across
runs and resumes.

% ======================================================================
\section{Optimizer Reference}
\label{sec:optimizers}

Each subsection gives the algorithm summary, a baseline configuration block
(directly usable in an experiment JSON), the meaning and valid range of every
option, and the issues one can run into.

% ----------------------------------------------------------------------
\subsection{\texttt{skopt\_gp}: Bayesian Optimization (GP surrogate)}
\label{sec:opt-bo}

A Gaussian-process surrogate over Expected Improvement, provided by
\texttt{scikit-optimize}.  The run starts with \texttt{n\_initial} Sobol'
quasi-random samples to build the initial surrogate, then switches to
model-guided proposals.  Naturally handles bounds and multi-modal
landscapes; the surrogate refit after every \texttt{tell()} makes it the
computationally heaviest optimizer per proposal (irrelevant when each
evaluation is a CFD run).

\begin{lstlisting}
"optimizer": {
    "kind": "skopt_gp",
    "stopping_criteria": { "n_calls": 20, "epsilon": 1e-4, "window": 3 },
    "kind_specific_options": {
        "n_initial": 8,
        "random_state": 42
    }
}
\end{lstlisting}

\begin{center}
\begin{tabular}{@{}llp{8.5cm}@{}}
\toprule
Option & Default & Description \\
\midrule
\texttt{n\_initial} & 8 & Sobol' samples before the GP takes over.
Must be $\geq 1$.  Rule of thumb: $8 \times D$ for a $D$-dimensional
problem when the budget allows. \\
\texttt{random\_state} & 42 & Seeds the Sobol' sampler \emph{and} the GP
hyperparameter-fit restarts, making runs fully reproducible. \\
\bottomrule
\end{tabular}
\end{center}

\paragraph{Possible issues.}
\begin{itemize}
  \item If \texttt{n\_initial} is close to \texttt{n\_calls}, the run is
        effectively quasi-random sampling with no model-guided phase.
  \item With a very small \texttt{n\_initial} the first GP fits rest on very
        few points and early model-guided proposals can be poor; the run
        still recovers, but wastes budget.
\end{itemize}

% ----------------------------------------------------------------------
\subsection{\texttt{nelder\_mead}: Nelder--Mead simplex}
\label{sec:opt-nm}

A derivative-free simplex method implemented as an explicit ask/tell state
machine.  Phase 1 evaluates the $D{+}1$ startup simplex points one at a
time; afterwards every ask/tell cycle performs one simplex operation
(reflect, expand, contract, or one shrink step).  The startup simplex is
built around the canonical GEKO defaults
($\csep{=}1.75$, $\cnw{=}0.5$, $\cmix{=}0.0$, $\cjet{=}1.0$,
$\ccorner{=}1.0$, $\cturb{=}2.0$; only the coefficients actually being
tuned are used) with a uniform offset of $0.1$ per dimension, clipped into
the bounds.  No \texttt{n\_initial} is needed.

\begin{lstlisting}
"optimizer": {
    "kind": "nelder_mead",
    "stopping_criteria": { "n_calls": 20, "epsilon": 1e-5, "window": 3 },
    "kind_specific_options": {
        "alpha": 0.8,
        "gamma": 1.5,
        "rho": 0.5,
        "sigma": 0.5
    }
}
\end{lstlisting}

\begin{center}
\begin{tabular}{@{}llp{8.5cm}@{}}
\toprule
Option & Default & Description \\
\midrule
\texttt{alpha} & 0.8 & Reflection coefficient. Must be $>0$.  Deliberately
below the textbook $1.0$: less aggressive reflection reduces overshoot in
the smooth, bowl-like CFD cost surface. \\
\texttt{gamma} & 1.5 & Expansion coefficient. Must be $>1$ (textbook
$2.0$; same rationale). \\
\texttt{rho} & 0.5 & Contraction coefficient. Must be in $(0,1)$. \\
\texttt{sigma} & 0.5 & Shrink coefficient. Must be in $(0,1)$. \\
\bottomrule
\end{tabular}
\end{center}

\paragraph{Possible issues.}
\begin{itemize}
  \item \textbf{Local method:} NM follows the local descent direction from
        the GEKO defaults; on a multi-modal cost surface it converges to the
        nearest minimum, not necessarily the global one
        (Section~\ref{sec:res-nm} shows exactly this).
  \item \textbf{Out-of-bounds proposals} are evaluated as-is
        (Section~\ref{sec:bounds}).
  \item In 1D the simplex has only two vertices, so ``second worst'' equals
        ``best'' and the plain accept-reflection branch is unreachable; every
        reflection triggers either an expansion attempt or a contraction.
        Inherent to 1D NM and harmless.
\end{itemize}

% ----------------------------------------------------------------------
\subsection{\texttt{finite\_differences}: forward-difference gradient descent}
\label{sec:opt-fd}

Classic gradient descent with a numerically estimated gradient.  Starting
from the GEKO defaults, each cycle costs $D{+}1$ evaluations: $D$ forward
probes ($\text{base} + \delta_i e_i$, one per dimension) to estimate the
gradient, then one descent step $\text{base} - \eta \nabla f$.  The stepped
point becomes the base for the next cycle unconditionally.  The probe
magnitude is
$\delta_i = \mathrm{clip}\bigl((\mathrm{high}_i - \mathrm{low}_i)\cdot
\texttt{step\_size},\ \texttt{min\_step},\ \texttt{max\_step}\bigr)$.

\begin{lstlisting}
"optimizer": {
    "kind": "finite_differences",
    "stopping_criteria": { "n_calls": 20, "epsilon": 1e-4, "window": 3 },
    "kind_specific_options": {
        "step_size": 0.05,
        "learning_rate": 0.2,
        "min_step": 1e-5,
        "max_step": 0.25
    }
}
\end{lstlisting}

\begin{center}
\begin{tabular}{@{}llp{8.5cm}@{}}
\toprule
Option & Default & Description \\
\midrule
\texttt{step\_size} & 0.05 & Probe perturbation as a fraction of each
parameter's range. Must be $>0$. \\
\texttt{learning\_rate} & 0.2 & Gradient-descent step size $\eta$. Must be
$>0$.  See the calibration note below. \\
\texttt{min\_step} & $10^{-5}$ & Minimum absolute probe perturbation. Must
be $>0$. \\
\texttt{max\_step} & 0.25 & Maximum absolute probe perturbation. Must be
$\geq$ \texttt{min\_step}. \\
\bottomrule
\end{tabular}
\end{center}

\paragraph{Possible issues.}
\begin{itemize}
  \item \textbf{Learning-rate calibration is the make-or-break setting.}
        The appropriate value is $\eta \approx 1/L$, where $L$ is the
        Lipschitz constant of the gradient (the maximum curvature) of the
        cost surface.  A too-large $\eta$ overshoots and produces a zig-zag
        walk that never settles.  The library default of $0.2$ is a generic
        value; the synthetic benchmark uses $0.015$
        (Section~\ref{sec:benchmark}) and the real CFD runs used $0.075$
        standalone and $0.045$ in the BO$\to$FD exploitation phase.
  \item \textbf{The step series is not monotonic}: the walk simply continues
        from wherever the last step landed, so never assume the last iterate
        is the best; report best-so-far.
  \item \textbf{Forward-difference bias:} near a smooth optimum the walk
        settles into a small limit cycle offset roughly $\delta/2$ from the
        true optimum.  Set \texttt{epsilon} on real runs so the experiment
        terminates instead of oscillating through the budget.
  \item \textbf{Probes consume budget:} of $N$ evaluations only about
        $N/(D{+}1)$ are actual descent steps; the rest are internal gradient
        probes.
  \item \textbf{Out-of-bounds proposals} are evaluated as-is
        (Section~\ref{sec:bounds}).
\end{itemize}

% ----------------------------------------------------------------------
\subsection{\texttt{pso}: Particle Swarm Optimization}
\label{sec:opt-pso}

A population method: \texttt{n\_particles} particles move under the standard
velocity update
\[
v \leftarrow w\,v
  + c_1 r_1 \,(\mathbf{x}_{pbest} - \mathbf{x})
  + c_2 r_2 \,(\mathbf{x}_{gbest} - \mathbf{x}),
\qquad r_1, r_2 \sim \mathcal{U}(0,1),
\]
with the inertia weight $w$ decaying linearly from $w_{start}$ to $w_{end}$
over the run.  Particle evaluations are serialized through the ask/tell
loop; one swarm iteration completes after \texttt{n\_particles}
evaluations.  Initial positions are drawn from a Sobol' sequence (seeded the
same way as the BO optimizers, so PSO and BO runs with equal
\texttt{random\_state} are comparable); initial velocities are zero, so the
first move is driven entirely by the cognitive/social attraction terms.
Velocity is clamped component-wise to $\pm v_{max}$
($= \texttt{v\_max\_frac} \times$ range); bounds are handled by absorption
(Section~\ref{sec:bounds}).

The number of swarm iterations is \emph{derived from the budget}, not
configured: $\texttt{max\_iter} = \texttt{n\_calls} / \texttt{n\_particles}
- 1$ (the init sweep costs \texttt{n\_particles} evaluations, each iteration
costs \texttt{n\_particles} more), and the inertia decay exactly spans those
iterations.

\begin{lstlisting}
"optimizer": {
    "kind": "pso",
    "stopping_criteria": { "n_calls": 200, "epsilon": 1e-4, "window": 5 },
    "kind_specific_options": {
        "n_particles": 10,
        "w_start": 0.9,
        "w_end": 0.4,
        "c1": 1.5,
        "c2": 1.5,
        "v_max_frac": 0.2,
        "random_state": 42
    }
}
\end{lstlisting}

\begin{center}
\begin{tabular}{@{}llp{8.5cm}@{}}
\toprule
Option & Default & Description \\
\midrule
\texttt{n\_particles} & 10 & Swarm size. Must be $\geq 2$. \\
\texttt{w\_start} & 0.9 & Initial inertia weight. Requires
$0 < w_{end} \leq w_{start} \leq 1$. \\
\texttt{w\_end} & 0.4 & Final inertia weight (reached at
\texttt{max\_iter}). \\
\texttt{c1} & 1.5 & Cognitive coefficient (pull toward personal best).
Must be $>0$. \\
\texttt{c2} & 1.5 & Social coefficient (pull toward global best).
Must be $>0$. \\
\texttt{v\_max\_frac} & 0.2 & Velocity cap per dimension as a fraction of
that dimension's range, in $(0,1]$.  Applied every iteration, not just at
initialization. \\
\texttt{random\_state} & 42 & Seeds the Sobol' initial positions and the
per-iteration $r_1$/$r_2$ draws. \\
\bottomrule
\end{tabular}
\end{center}

\paragraph{Possible issues.}
\begin{itemize}
  \item \textbf{Budget divisibility:} \texttt{test\_config()} raises unless
        \texttt{n\_calls} is divisible by \texttt{n\_particles} (so the
        budget divides into whole swarm iterations) and
        $\texttt{n\_calls} \geq 2 \times \texttt{n\_particles}$ (init sweep
        plus at least one swarm iteration).
  \item \textbf{Particles vs.\ iterations trade-off:} with a fixed budget,
        more particles means broader coverage per iteration but fewer
        iterations to converge.  The real-CFD sweep (see the companion
        comparison report) found small swarms (3--5 particles) preferable
        at tight budgets.
  \item The epsilon check advances only once per swarm iteration, so with
        large \texttt{n\_particles} the early-stopping machinery needs many
        evaluations before it can fire ($2\times\texttt{window}$
        iterations).
\end{itemize}

% ----------------------------------------------------------------------
\subsection{\texttt{hybrid\_nm\_bayes}: Nelder--Mead $\to$ Bayesian}
\label{sec:opt-nmbo}

Warm-start hybrid.  Phase 1: the first \texttt{n\_initial} evaluations run
Nelder--Mead (self-initialized from the GEKO defaults, exactly as in
Section~\ref{sec:opt-nm}).  Phase 2: the remaining
$\texttt{n\_calls} - \texttt{n\_initial}$ evaluations run a GP that was
warm-started with the NM history (the BO optimizer is constructed with zero
initial points of its own and receives every in-bounds NM result via
\texttt{tell()}).

\begin{lstlisting}
"optimizer": {
    "kind": "hybrid_nm_bayes",
    "stopping_criteria": { "n_calls": 22, "epsilon": 1e-4, "window": 3 },
    "kind_specific_options": {
        "n_initial": 10,
        "bo_options": { "bayesian_kind": "GP", "random_state": 42 },
        "nm_options": { "alpha": 0.8, "gamma": 1.5, "rho": 0.5, "sigma": 0.5 }
    }
}
\end{lstlisting}

\begin{center}
\begin{tabular}{@{}llp{8.2cm}@{}}
\toprule
Option & Default & Description \\
\midrule
\texttt{n\_initial} & 8 & Number of NM evaluations before the switch.
Must be $\geq 1$. \\
\texttt{bo\_options.bayesian\_kind} & \texttt{"GP"} & Phase-2 surrogate:
\texttt{"GP"}, \texttt{"RF"}, \texttt{"ET"}, or \texttt{"GBRT"}.
\texttt{"GP"} resolves to the same explicitly built estimator as
\texttt{skopt\_gp}. \\
\texttt{bo\_options.random\_state} & 42 & Seed for the Bayesian phase. \\
\texttt{nm\_options} & (see \ref{sec:opt-nm}) & Same options as standalone
\texttt{nelder\_mead}. \\
\bottomrule
\end{tabular}
\end{center}

\paragraph{Possible issues.}
\begin{itemize}
  \item \textbf{Out-of-bounds warm-up points are dropped from the BO warm
        start:} NM may propose points outside the bounds
        (Section~\ref{sec:bounds}), but \texttt{skopt} rejects points
        outside its space, so those evaluations are silently excluded from
        the GP's training history.  In the pathological extreme (zero
        in-bounds warm-up points) the first BO-phase \texttt{ask()} raises
        a \texttt{RuntimeError}; this cannot happen under normal conditions
        because the warm-up starts in-bounds at the GEKO defaults, and is
        deliberately left unhandled.
  \item \textbf{Inherited local bias:} the GP enters its phase trained
        almost exclusively on points near wherever NM went.  If NM descended
        into a local minimum, the warm-started BO phase needs several
        exploratory evaluations to look elsewhere; with a short Phase 2
        it may not escape at all (Section~\ref{sec:res-nmbo}).
\end{itemize}

% ----------------------------------------------------------------------
\subsection{\texttt{hybrid\_fd\_bayes}: Finite Differences $\to$ Bayesian}
\label{sec:opt-fdbo}

Warm-start hybrid, identical in structure to \texttt{hybrid\_nm\_bayes} but
with finite-difference gradient descent as the Phase-1 method.

\begin{lstlisting}
"optimizer": {
    "kind": "hybrid_fd_bayes",
    "stopping_criteria": { "n_calls": 60, "epsilon": 1e-4, "window": 3 },
    "kind_specific_options": {
        "n_initial": 20,
        "bo_options": { "bayesian_kind": "GP", "random_state": 42 },
        "fd_options": {
            "step_size": 0.05, "learning_rate": 0.2,
            "min_step": 1e-5, "max_step": 0.25
        }
    }
}
\end{lstlisting}

\texttt{bo\_options} as in Section~\ref{sec:opt-nmbo};
\texttt{fd\_options} takes the same options as standalone
\texttt{finite\_differences}; \texttt{n\_initial} (default 8, $\geq 1$) is
the number of FD evaluations before the switch.

\paragraph{Possible issues.}
\begin{itemize}
  \item All issues of standalone FD apply during Phase 1 (learning-rate
        calibration, probes consuming budget).
  \item Choose \texttt{n\_initial} so the FD phase ends at a cycle boundary:
        Phase 1 begins with one base evaluation and then runs cycles of
        $D{+}1$ evaluations, so $\texttt{n\_initial} = 1 + k\,(D{+}1)$ uses
        the FD budget without a dangling half-finished gradient estimate.
\end{itemize}

% ----------------------------------------------------------------------
\subsection{\texttt{hybrid\_bayes\_nm}: Bayesian $\to$ Nelder--Mead}
\label{sec:opt-bonm}

Refinement hybrid.  Phase 1: the first \texttt{n\_initial} evaluations run
BO (Sobol' sampling for \texttt{n\_initial\_sobol} points, then GP
proposals).  Phase 2: Nelder--Mead refinement, with its startup simplex
re-centred on the best point found by BO (instead of the GEKO defaults),
using the same uniform $0.1$ offsets scaled by \texttt{simplex\_scale}.

\begin{lstlisting}
"optimizer": {
    "kind": "hybrid_bayes_nm",
    "stopping_criteria": { "n_calls": 30, "epsilon": 1e-4, "window": 3 },
    "kind_specific_options": {
        "n_initial": 10,
        "bo_options": {
            "n_initial_sobol": 5, "bayesian_kind": "GP", "random_state": 42
        },
        "nm_options": {
            "alpha": 0.8, "gamma": 1.5, "rho": 0.5, "sigma": 0.5,
            "simplex_scale": 0.6
        }
    }
}
\end{lstlisting}

\begin{center}
\begin{tabular}{@{}llp{7.6cm}@{}}
\toprule
Option & Default & Description \\
\midrule
\texttt{n\_initial} & 8 & Number of BO evaluations before the switch.
Must be $\geq 1$. \\
\texttt{bo\_options.n\_initial\_sobol} & $\min(5, \texttt{n\_initial})$ &
Sobol' samples within the BO phase before the GP takes over. \\
\texttt{bo\_options.bayesian\_kind} & \texttt{"GP"} & As in
Section~\ref{sec:opt-nmbo}. \\
\texttt{bo\_options.random\_state} & 42 & Seed for Sobol' and the GP. \\
\texttt{nm\_options.simplex\_scale} & 1.0 & Scale factor for the NM startup
simplex built around the BO best point; values $<1$ tighten local
exploitation (e.g.\ $0.6$ shrinks the offsets to $\pm 0.06$).  Ignored by
the standalone \texttt{nelder\_mead} kind. \\
\texttt{nm\_options} (rest) & (see \ref{sec:opt-nm}) & Same options as
standalone \texttt{nelder\_mead}. \\
\bottomrule
\end{tabular}
\end{center}

\paragraph{Possible issues.}
\begin{itemize}
  \item If the BO phase has not yet localized a good region when
        \texttt{n\_initial} runs out, NM refines around a mediocre point;
        the refinement is only as good as the hand-off.  Give BO enough
        budget to finish its Sobol' phase plus a meaningful number of
        model-guided proposals.
\end{itemize}

% ----------------------------------------------------------------------
\subsection{\texttt{hybrid\_bayes\_fd}: Bayesian $\to$ Finite Differences}
\label{sec:opt-bofd}

Refinement hybrid.  Phase 1 as in \texttt{hybrid\_bayes\_nm}.  Phase 2:
finite-difference descent starting from the BO best point, which is injected
directly as the FD base \emph{without re-evaluating it}; the FD phase
begins immediately with a gradient probe rather than wasting an evaluation
on the already-known base.

\begin{lstlisting}
"optimizer": {
    "kind": "hybrid_bayes_fd",
    "stopping_criteria": { "n_calls": 30, "epsilon": 1e-4, "window": 3 },
    "kind_specific_options": {
        "n_initial": 10,
        "bo_options": {
            "n_initial_sobol": 5, "bayesian_kind": "GP", "random_state": 42
        },
        "fd_options": {
            "step_size": 0.03, "learning_rate": 0.009,
            "min_step": 1e-5, "max_step": 0.25
        }
    }
}
\end{lstlisting}

\texttt{bo\_options} as in Section~\ref{sec:opt-bonm}; \texttt{fd\_options}
takes the same options as standalone \texttt{finite\_differences}.

\paragraph{Possible issues.}
\begin{itemize}
  \item \textbf{Tighten the FD settings:} because BO has already localized a
        promising region, the FD phase should use a smaller
        \texttt{step\_size} and \texttt{learning\_rate} than a standalone FD
        run; roughly $60\%$ of the standalone values is a reasonable
        starting point (e.g.\ $0.03$ instead of $0.05$).  A more local
        gradient estimate and a smaller step avoid overshooting the region
        BO identified.
  \item The FD phase budget ($\texttt{n\_calls} - \texttt{n\_initial}$) is
        spent in cycles of $D{+}1$ evaluations starting directly with a
        probe; choosing it as a multiple of $D{+}1$ avoids a dangling
        half-finished gradient estimate.
  \item All standalone-FD caveats apply (non-monotonic step series,
        forward-difference bias).
\end{itemize}

% ======================================================================
\section{Synthetic Benchmark: Testing Routine}
\label{sec:benchmark}

Before being applied to real CFD cases, all eight optimizers are exercised
on cheap analytic test functions by the scripts in
\texttt{optimizer\_visualization/}: \texttt{\_benchmark\_core.py} holds the
test functions, optimizer configurations, and plotting helpers, and
\texttt{benchmark\_individual.py} produces one standalone figure per
optimizer (a side-by-side pair of the 1-D trajectory and the 2-D contour
trajectory, written to \texttt{plots/individual/}).

\subsection{Test functions}
\label{sec:testfns}

Both test functions are deliberately \emph{bimodal}: they possess a local
and a global minimum, with the GEKO default starting point placed near a
local \emph{maximum} between them, a trap for purely local methods.

\paragraph{1-D.} On $\csep \in [0.5, 3.5]$:
\begin{equation}
f_{1D}(x) = -2.5x \,\sin(2.5x).
\label{eq:f1d}
\end{equation}
It has a local minimum at $x \approx 0.811$ ($f \approx -1.820$) and the
global minimum at $x \approx 3.191$ ($f \approx -7.917$).  The GEKO default
$x = 1.75$ sits at $f \approx +4.13$, near the local maximum between the two
minima: a gradient follower started there descends into the \emph{local}
minimum.

\paragraph{2-D.} On $(\csep, \cnw) \in [0.5, 3.5] \times [0.1, 0.9]$:
\begin{equation}
f_{2D}(x_1, x_2) = f_{1D}(x_1) + 3\,\lvert x_2 - 0.5 \rvert.
\label{eq:f2d}
\end{equation}
The $x_1$ direction reuses the 1-D objective and keeps its local/global
structure; the $x_2$ term adds a V-shaped valley with a \emph{non-smooth}
floor at $x_2 = 0.5$ (intentional: none of the optimizers use analytic
gradients, and the kink probes how the numerical-gradient and simplex
methods cope with non-smoothness).  Global minimum at $(\approx 3.191,\,
0.5)$ with $f \approx -7.917$; local minimum at $(\approx 0.811,\, 0.5)$.
The GEKO default $(1.75, 0.5)$ again sits near the 1-D local maximum, on the
$x_2$ valley floor.

\subsection{Harness}

Every optimizer is built through the production
\texttt{build\_optimizer()} path from an \texttt{OptimizerSection}; the
benchmark exercises the exact code used in real experiments.  The loop is
the plain ask/tell alternation with a fixed budget of $N = 24$ evaluations
in 1-D and $N = 48$ in 2-D, identical for all optimizers.  Proposals are
evaluated exactly where suggested, without clipping, so the NM/FD boundary
behavior of Section~\ref{sec:bounds} is visible in the plots rather than
masked.

\paragraph{Plot elements.} In every figure: the red star is the true global
minimum, the gold star the best point found; dot color encodes evaluation
order (dark purple first, bright yellow last); for hybrids, circles are
Phase-1 and triangles Phase-2 evaluations; for PSO each particle keeps one
fixed color so its path is traceable.  For the FD-based optimizers the
internal gradient \emph{probes} are hidden from the plot (they are
bookkeeping, not decisions), but the ``$N = \ldots$ function evals''
annotation still counts them: in 1-D, $N=24$ evaluations produce 12 visible
dots (1 base $+$ 11 steps), in 2-D $N=48$ produce 16 (1 base $+$ 15 steps).

\subsection{Optimizer initialization}
\label{sec:bench-init}

Table~\ref{tab:bench-config} lists the exact configuration of every
optimizer in the benchmark (anything not listed uses the library defaults of
Section~\ref{sec:optimizers}).  All optimizers use
\texttt{random\_state=42}, and every optimizer with an initialization phase
uses $\texttt{n\_initial}=9$ in 1-D and $\texttt{n\_initial}=15$ in 2-D.
The BO-first hybrids allocate a longer Phase 1 (17 evaluations in 1-D, 32
in 2-D), which leaves room for both the Sobol' initialization and a
model-guided stretch before the hand-off; their refinement phases are
tightened as recommended in
Sections~\ref{sec:opt-bonm}--\ref{sec:opt-bofd}.

The FD learning rate is calibrated to the test function: the maximum
curvature of $f_{1D}$ on the domain is
$L = \max \lvert f_{1D}'' \rvert \approx 52$, so the benchmark uses
$\eta = 0.015 < 1/L \approx 0.019$.  The $x_2$ direction of $f_{2D}$ is
piecewise linear (slope $\pm 3$), so the same rate is safe there.

\begin{table}[htbp]
\centering
\small
\begin{tabular}{@{}llll@{}}
\toprule
Optimizer & \texttt{kind} & 1-D configuration ($N{=}24$) & 2-D configuration ($N{=}48$) \\
\midrule
Bayesian Opt. & \texttt{skopt\_gp} &
  \texttt{n\_initial=9} & \texttt{n\_initial=15} \\
Nelder--Mead & \texttt{nelder\_mead} &
  defaults & defaults \\
Finite Differences & \texttt{finite\_differences} &
  \texttt{step=0.05, lr=0.015} & \texttt{step=0.05, lr=0.015} \\
Particle Swarm & \texttt{pso} &
  \texttt{4 particles} & \texttt{4 particles} \\
 & & (5 swarm iterations) & (11 swarm iterations) \\
Hybrid NM$\to$BO & \texttt{hybrid\_nm\_bayes} &
  \texttt{n\_initial=9} & \texttt{n\_initial=15} \\
Hybrid FD$\to$BO & \texttt{hybrid\_fd\_bayes} &
  \texttt{n\_initial=9}, & \texttt{n\_initial=15}, \\
 & & \texttt{fd: step=0.05, lr=0.015} & \texttt{fd: step=0.05, lr=0.015} \\
Hybrid BO$\to$NM & \texttt{hybrid\_bayes\_nm} &
  \texttt{n\_initial=17, sobol=9}, & \texttt{n\_initial=32, sobol=5}, \\
 & & \texttt{nm: simplex\_scale=0.6} & \texttt{nm: simplex\_scale=0.6} \\
Hybrid BO$\to$FD & \texttt{hybrid\_bayes\_fd} &
  \texttt{n\_initial=17, sobol=9}, & \texttt{n\_initial=32, sobol=15}, \\
 & & \texttt{fd: step=0.03, lr=0.009} & \texttt{fd: step=0.03, lr=0.009} \\
\bottomrule
\end{tabular}
\caption{Benchmark configurations (\texttt{step} = \texttt{step\_size},
\texttt{lr} = \texttt{learning\_rate}, \texttt{sobol} =
\texttt{n\_initial\_sobol}).  All optimizers use
\texttt{random\_state=42}.  PSO's swarm-iteration counts are derived from
the budget: $24/4 - 1 = 5$ and $48/4 - 1 = 11$.}
\label{tab:bench-config}
\end{table}

% ======================================================================
\section{Synthetic Benchmark: Results}
\label{sec:benchmark-results}

For reference in the discussions below: the true global minimum is
$f^\ast = -7.917$ (at $x \approx 3.191$ in 1-D, at $(\approx 3.191,\, 0.5)$
in 2-D) and the local minimum is $f \approx -1.820$.

% ----------------------------------------------------------------------
\subsection{Bayesian Optimization}
\label{sec:res-bo}

\begin{center}
\includegraphics[width=\textwidth]{01_bayesian_opt.png}
\end{center}

Works as intended in both dimensions.  The Sobol' samples (9 in 1-D, 15 in
2-D) cover both basins; the model-guided phase then concentrates its
remaining evaluations in the global basin and lands essentially on the true
optimum: best $x = 3.191$, $f = -7.9167$ in 1-D and best $(3.19,\ 0.50)$,
$f = -7.9162$ in 2-D.  The local-minimum trap is avoided entirely.

% ----------------------------------------------------------------------
\subsection{Nelder--Mead}
\label{sec:res-nm}

\begin{center}
\includegraphics[width=\textwidth]{02_nelder_mead.png}
\end{center}

Works exactly as a local method should, which here means converging to
the \emph{local} minimum.  Starting from the GEKO-default trap point, the
simplex walks downhill and contracts tightly onto $x = 0.812$,
$f = -1.8197$ in 1-D (the local optimum to four decimals) and
$(0.81,\ 0.50)$, $f = -1.8193$ in 2-D.  The global basin is never sampled;
this is the expected and intended failure mode of a purely local method on
a bimodal landscape, not an implementation defect.

% ----------------------------------------------------------------------
\subsection{Finite Differences}
\label{sec:res-fd}

\begin{center}
\includegraphics[width=\textwidth]{03_finite_differences.png}
\end{center}

Works as intended: a clean descent from the trap point into the local basin
(only bases and gradient steps are plotted; the probes are hidden, so 12 of
24 and 16 of 48 evaluations are visible).  In 1-D the walk settles
essentially on the local optimum, $x = 0.809$, $f = -1.8197$; the tiny
remaining offset from $x = 0.811$ is consistent with the forward-difference
bias of Section~\ref{sec:opt-fd}.  In 2-D it reaches $(0.79,\ 0.50)$,
$f = -1.8157$, handling the non-smooth $x_2$ ridge without oscillation at
this learning rate.  Like NM, it cannot see the global basin.

% ----------------------------------------------------------------------
\subsection{Particle Swarm Optimization}
\label{sec:res-pso}

\begin{center}
\includegraphics[width=\textwidth]{04_particle_swarm.png}
\end{center}

Works: the Sobol'-initialized swarm covers both basins, the global best
immediately anchors in the global basin, and the swarm contracts onto it.
Best $x = 3.169$, $f = -7.9038$ in 1-D and $(3.19,\ 0.50)$, $f = -7.9153$
in 2-D, marginally above the true optimum, as expected from only 5 (1-D)
resp.\ 11 (2-D) swarm iterations with four particles.  The per-particle
coloring shows all four particles ending in the global basin.

% ----------------------------------------------------------------------
\subsection{Hybrid NM $\to$ BO}
\label{sec:res-nmbo}

\begin{center}
\includegraphics[width=\textwidth]{05_hybrid_nm_bo.png}
\end{center}

Illustrates the warm-start hybrid's inherited local bias
(Section~\ref{sec:opt-nmbo}).  In 1-D the NM phase descends into the local
basin and the warm-started GP, trained only on those local points, spends
its fifteen evaluations refining there: best $x = 0.812$, $f = -1.8197$,
the local minimum.  In 2-D the longer BO phase does break out: it discovers
the global basin and reaches $(3.18,\ 0.44)$, $f = -7.7394$, far below the
local minimum but slightly off the $x_2 = 0.5$ ridge.

A second observation stands out in the 2-D panel: once the BO phase finds
the global basin, its evaluations concentrate in a narrow column around the
best point found instead of spreading over the still-unexplored parts of
the domain.  The BO phase behaves strongly exploitative rather than
explorative here.  This is presumably a phenomenon of the underlying
Gaussian process (its fitted hyperparameters and the resulting acquisition
landscape) and needs to be studied further.

% ----------------------------------------------------------------------
\subsection{Hybrid FD $\to$ BO}
\label{sec:res-fdbo}

\begin{center}
\includegraphics[width=\textwidth]{06_hybrid_fd_bo.png}
\end{center}

Same pattern as NM$\,\to\,$BO.  In 1-D the FD warm-up walks toward the
local basin and the warm-started BO phase stays there: best $x = 0.812$,
$f = -1.8197$.  In 2-D the BO phase escapes to the global basin, reaching
$f = -7.7786$ close to the optimum but slightly off the $x_2 = 0.5$ ridge.
The same strongly exploitative Phase-2 behavior as in
Section~\ref{sec:res-nmbo} is visible here: the BO evaluations stack up in
a tight column around the best point rather than exploring; see the remark
there.

% ----------------------------------------------------------------------
\subsection{Hybrid BO $\to$ NM}
\label{sec:res-bonm}

\begin{center}
\includegraphics[width=\textwidth]{07_hybrid_bo_nm.png}
\end{center}

Works, essentially perfectly.  The BO exploration phase finds the global
basin; the NM refinement phase, re-centred on the BO best with its
tightened simplex (\texttt{simplex\_scale=0.6}), converges onto the
optimum: best $x = 3.191$, $f = -7.9167$ in 1-D and $(3.19,\ 0.50)$,
$f = -7.9163$ in 2-D.  The plots show the Phase-2 evaluations clustered
tightly around the true optimum, the intended
exploration-then-exploitation division of labor.

% ----------------------------------------------------------------------
\subsection{Hybrid BO $\to$ FD}
\label{sec:res-bofd}

\begin{center}
\includegraphics[width=\textwidth]{08_hybrid_bo_fd.png}
\end{center}

Works.  In 1-D the configuration of the BO phase is identical to
BO$\,\to\,$NM (same budget, Sobol' count, and seed), so the hand-off point
is the same and so is the best found: $x = 3.191$, $f = -7.9167$, already
located during the BO phase.  In 2-D the tightened FD refinement
(\texttt{step\_size=0.03}, \texttt{learning\_rate=0.009}) descends within
the global basin to $(3.20,\ 0.50)$, $f = -7.9039$, slightly above
BO$\,\to\,$NM's result.  That gap is consistent with FD's probe overhead
($D{+}1$ evaluations per descent step) and forward-difference bias, and
matches the real-CFD comparison, where the FD-refined hybrid also trailed
the NM-refined one.

% ----------------------------------------------------------------------
\subsection{Summary}

\begin{table}[htbp]
\centering
\begin{tabular}{@{}lllll@{}}
\toprule
 & \multicolumn{2}{c}{1-D ($f^\ast = -7.917$)} &
   \multicolumn{2}{c}{2-D ($f^\ast = -7.917$)} \\
Optimizer & best $f$ & minimum found & best $f$ & minimum found \\
\midrule
Bayesian Opt.      & $-7.9167$ & global & $-7.9162$ & global \\
Nelder--Mead       & $-1.8197$ & local  & $-1.8193$ & local  \\
Finite Differences & $-1.8197$ & local  & $-1.8157$ & local \\
Particle Swarm     & $-7.9038$ & global & $-7.9153$ & global \\
Hybrid NM$\to$BO   & $-1.8197$ & local  & $-7.7394$ & global basin, off-ridge \\
Hybrid FD$\to$BO   & $-1.8197$ & local  & $-7.7786$ & global basin, off-ridge \\
Hybrid BO$\to$NM   & $-7.9167$ & global & $-7.9163$ & global \\
Hybrid BO$\to$FD   & $-7.9167$ & global & $-7.9039$ & global \\
\bottomrule
\end{tabular}
\caption{Best objective value found by each optimizer within the shared
budget (24 evaluations in 1-D, 48 in 2-D), read from the benchmark figures.}
\label{tab:bench-summary}
\end{table}

The benchmark confirms the intended division: every optimizer behaves
according to its design.  Methods with a global-exploration component (BO,
PSO, and the BO-first hybrids) find the global minimum; purely local
methods (NM, FD) converge into the local-minimum trap by construction; and
the warm-start hybrids land in between, escaping the trap only when their
BO phase is long enough.  On the real CFD objective, which is much closer
to unimodal over the tested coefficient ranges, this trap does not
materialize, and all strategies converge to comparable optima (see the
companion comparison report).

\end{document}
